\chapter{방정식은 어떤 대칭을 갖는가}
\label{chap:volume3-motivation}

다항식
\[
  f(X)=a_nX^n+\cdots+a_1X+a_0\in K[X]
\]
의 근을 구한다는 말은 단순히 복소수 근사값을 계산한다는 뜻만은 아니다. 갈루아 이론의 질문은 더 구조적이다.

\begin{itemize}
  \item 근들을 서로 바꾸어도 계수체 $K$에서 보이는 모든 관계가 그대로 유지되는가?
  \item 가능한 근의 순열 가운데 실제 체자동동형에서 오는 것은 무엇인가?
  \item 그 순열들의 군은 방정식의 인수분해, 중간체와 근호해법을 어떻게 기록하는가?
\end{itemize}

분리다항식 $f$의 분해체를 $L$이라 하고 근들을
\[
  \alpha_1,\dots,\alpha_n
\]
이라 하자. 각 $K$-자동동형은 계수를 고정하므로
\[
  f(\sigma(\alpha_i))=\sigma(f(\alpha_i))=0
\]
이다. 따라서 갈루아군은 근들을 순열하고
\[
  \Gal(L/K)\hookrightarrow S_n
\]
로 들어간다. 그러나 모든 순열이 허용되는 것은 아니다. 예를 들어 두 근 사이에
\[
  \alpha_1+\alpha_2=0
\]
이라는 관계가 있으면 자동동형은 그 관계도 보존해야 한다.

\begin{example}[같은 차수, 다른 대칭]
다음 두 사차다항식은 모두 $\Q$ 위에서 분리되지만 분해체의 대칭은 다르다.
\[
  X^4-2,
  \qquad
  X^4-4X^2+16.
\]
첫 번째 분해체의 갈루아군은 위수 $8$의 이면군과 동형이고, 두 번째는 원시 $12$차 단위근이 생성하는 원분체에서 클라인 네 군을 준다. 다항식의 차수만으로 갈루아군을 결정할 수 없다는 뜻이다.
\end{example}

\begin{keyidea}
방정식의 갈루아군을 계산하는 일은 다음 네 층을 연결하는 작업이다.
\[
  \boxed{
  \text{인수분해}
  \longleftrightarrow
  \text{근의 궤도}
  \longleftrightarrow
  \text{분해체의 차수}
  \longleftrightarrow
  \text{순열부분군}
  }
\]
\end{keyidea}

이 권의 첫 장에서는 이 연결을 이차, 삼차와 특별한 사차방정식에서 완전히 계산한다. 특히 삼차방정식에서는 판별식이 갈루아군 안에 홀순열이 존재하는지를 감지한다. 뒤의 장에서는 이 관찰을 일반 차수의 판별식과 종결식으로 확장한다.
